Analizzando l'incertezza sulle misure della velocità della luce $c$, si ritrovano due diverse categorie di errori che la generano:
\begin{itemize}
    \item errore sistematico, dovuto all'incertezza sulle misure $D$, $a$ ed $f_2$. Mediante il metodo della propagazione degli errori, si può calcolare $\sigma_{sis}$.
    Tale quantità corrisponde ad una incertezza presente in tutte le misure, che non può essere ridotta utilizzando metodi statistici.
    
    \item errore casuale, dovuto all'incertezza presente sulle misure $\Delta\delta$ e $\Delta\omega$, che variano di misura in misura. In particolare, l'incertezza che caratterizza la misura di $\Delta\delta$ è molto significativa, poiché con uno strumento dotato di precisione $0.01mm$ si misurano distanza comprese tra $0. 05mm $ e $1mm$ (ovviamente tale criticità risulta meno importante nel caso del turbo).
    L'incertezza sulla misura della velocità dovuta all'errore casuale può essere ridotta usando strumenti statistici, in particolare la deviazione standard della media, data dalla formula:
    \begin{equation}
         \sigma_{\bar{c}}=\sqrt{\frac{\sum_{i=1}^{n} (c_i-\bar{c})^2}{N(N-1)}}
    \end{equation}
\end{itemize}

L'incertezza totale sulla misura finale di $c$ è dunque dato da
\begin{equation}
    \sigma_c=\sqrt{\sigma_{\bar{c}}^2+\sigma_{sis}^2}
\end{equation}

In particolare si ricaveranno le misure di c per i tre differenti metodi di sperimentazione (rotazione dello specchio in senso orario, in senso antiorario e usando il turbo nei due sensi), in modo da assicurarsi che le misure in ogni set di dati siano ottenute con procedure e strumentazioni uguali.